%   terminology:
%     process: refers to execution in user space, or maybe struct proc &c
%     process memory: the lower part of the address space
%     process has one thread with two stacks (one for in kernel mode and one for
%     in user mode)
% talk a little about initial page table conditions:
%     paging not on, but virtual mostly mapped direct to physical,
%     which is what things look like when we turn paging on as well
%     since paging is turned on after we create first process.
%   mention why still have SEG_UCODE/SEG_UDATA?
%   do we ever really say what the low two bits of %cs do?
%     in particular their interaction with PTE_U
%   sidebar about why it is extern char[]

\chapter{Operating system organization}
\label{CH:FIRST}

对操作系统的一个关键要求是同时支持多个活动。例如，
人们可以使用 Chapter~\ref{CH:UNIX} 中的 {\tt fork} 和 {\tt exec} 系统调用
来启动一个编译器和一个文本编辑器作为进程。
操作系统必须在这些进程之间\indextext{time-share}（时分复用）
CPU 和内存等资源。
操作系统还必须安排这些进程之间的\indextext{isolation}（隔离）。
如果一个进程有 bug 并发生故障，
它不应该影响不相关的进程。
然而，完全的隔离又过于强硬，因为进程之间应该能够
有意地进行交互；管道就是一个例子。
因此，
一个操作系统必须满足三个要求：多路复用（multiplexing）、隔离（isolation）
和交互（interaction）。
本章概述了操作系统是如何组织
以实现这三个要求的。
事实证明有很多方法可以做到这一点，但本书重点关注
以\indextext{monolithic kernel}（宏内核）为中心的网络主流设计，这种设计被许多
Unix 操作系统所使用。
本章还概述了 
xv6 进程，它是 xv6 中的隔离单元。
Xv6 运行在一个\indextext{multi-core}\footnote{
本书所说的“多核”是指共享内存但并行执行的多个 CPU，每个 CPU 都有自己的一组寄存器。
本书有时会使用术语\indextext{multiprocessor}（多处理器）作为多核的同义词，
尽管多处理器也可以更具体地指具有几个不同处理器芯片的计算机。}RISC-V 微处理器上，
并且它的许多低级功能（例如其进程实现）是特定于 
RISC-V 的。
RISC-V 是一个 64 位的 CPU，而 xv6 是用 “LP64” C 语言编写的，
这意味着 C 语言中的 long (L) 和指针 (P) 
是 64 位的，但 {\tt int} 是 32 位的。
本书假设读者已经在某些架构上进行过
一些机器级别的编程，并且将在出现时引入 
RISC-V 特有的概念。
用户级 ISA~\cite{riscv:user} 和特权
架构~\cite{riscv:priv} 文档是完整的规范。
你也可以参考 
“The RISC-V Reader: An Open Architecture Atlas”~\cite{riscv}。
完整计算机中的 CPU 被支持硬件所包围，其中大部分
是以 I/O 接口的形式存在。
Xv6 是为 qemu 的 “-machine virt” 选项所模拟的支持硬件而编写的。
这包括 
RAM、包含引导代码的 ROM、到用户键盘/屏幕的串口连接以及用于存储的磁盘。
%% 
\section{Abstracting physical resources}
%% 

在遇到操作系统时，人们可能会问的第一个问题是为什么要
有它？
也就是说，人们可以将 Figure~\ref{fig:api} 中的系统调用
实现为一个库，应用程序与其进行链接。
在这种方案中，
每个应用程序甚至可以拥有一个根据其需求量身定制的自己的库。
应用程序可以直接与硬件资源交互，
并以最适合该应用程序的方式使用这些资源（例如，以实现高效或可预测的性能）。
一些用于嵌入式设备或实时系统的操作系统就是以这种方式组织的。
这种库方法的不利之处在于，如果有多个应用程序正在运行，这些应用程序必须表现良好。
例如，每个应用程序必须定期放弃 
CPU，以便其他应用程序可以运行。
如果所有应用程序都互相信任并且没有 bug，
这种\textit{cooperative}（协作式）的时分复用方案可能是可以的。
更典型的情况是应用程序之间互不信任，并且存在 bug，因此人们通常希望得到
比协作方案提供的更强的隔离。
为了实现强隔离，禁止应用程序
直接访问敏感的硬件资源，而是将资源抽象为服务，这是很有帮助的。
例如，Unix 应用程序仅通过文件系统的 
\lstinline{open}、\lstinline{read}、\lstinline{write} 和 \lstinline{close} 
系统调用与存储进行交互，
而不是直接对磁盘进行读写。
这为应用程序提供了路径名的便利，并且它允许
提供该接口的操作系统来管理磁盘。
即使不考虑隔离问题，
故意进行交互（或者只是希望互不干扰）的程序
很可能会发现文件系统是一个比直接使用磁盘更方便的抽象。
类似地，Unix 在进程之间透明地切换硬件 CPU，
根据需要保存和恢复寄存器状态，
这样应用程序就不需要意识到时分复用的存在。
这种透明性允许操作系统共享 
CPU，即使某些应用程序处于死循环中。
作为另一个例子，Unix 进程使用 \lstinline{exec} 
来构建它们的内存映像，而不是直接与物理内存交互。
这允许操作系统决定将进程放在内存中的什么位置；
如果内存紧张，操作系统甚至可以将进程的一些
数据存储在磁盘上。
\lstinline{exec} 还为
用户提供了使用文件系统来存储可执行程序映像的便利。
Unix 进程之间的许多交互形式都是通过文件描述符进行的。
文件描述符不仅抽象掉了许多细节（例如，
管道或文件中的数据存储在哪里），而且它们的设计方式
也简化了交互。
例如，如果管道中的一个应用程序退出或失败，内核
会自动为管道中的下一个进程生成一个文件结束（end-of-file）信号。
Figure~\ref{fig:api} 中的系统调用接口
经过精心设计，既提供了程序员的便利，又提供了
强隔离的可能性。
Unix 接口不是抽象资源的唯一方法，但它已被证明是一个好方法。
%% 
\section{User mode, supervisor mode, and system calls}
%% 

强隔离需要应用程序和操作系统之间有一条坚固的边界。
不应该允许应用程序干扰操作系统或其他程序的运行，即使该应用程序
有 bug 或是恶意的。
为了实现强隔离，操作系统必须安排好让应用程序无法修改（甚至无法读取）操作系统的内核数据结构和指令，并且
应用程序无法访问其他进程的内存。
CPU 为强隔离提供硬件支持。例如，RISC-V 有三个特权级别，它们限制了
代码可以做什么：
\indextext{machine mode}（机器模式）、
\indextext{supervisor mode}（主管模式）和
\indextext{user mode}（用户模式）。
在机器模式下执行的指令具有最高特权；
CPU 在机器模式下启动。
机器模式主要用于在引导期间设置计算机。
Xv6 在机器模式下短暂执行，
然后更改为主管模式。
在主管模式下，CPU 被允许执行\indextext{privileged instructions}（特权指令）：
例如，允许和禁止中断、读取和写入保存页表地址的寄存器等。
如果用户模式下的应用程序尝试执行
一条特权指令，那么 CPU 不会执行该指令，而是
“trap”（陷阱）到主管模式下的特殊代码中，该代码可以终止该应用程序。
Chapter~\ref{CH:UNIX} 中的 Figure~\ref{fig:os} 阐明了这种组织形式。一个应用程序只能
执行用户模式指令（例如，数字相加等），并被称为运行在\indextext{user space}（用户空间）中，
而主管模式下的软件也可以执行特权指令，并被称为运行在\indextext{kernel space}（内核空间）中。
运行在内核空间（或主管模式下）的软件被称为
\indextext{kernel}（内核）。
应用程序通过系统调用（例如 {\tt read}）与内核交互。
不允许应用程序直接调用内核函数或访问内核的内存。
RISC-V 为系统调用提供了 \indexcode{ecall} 指令；
它将 CPU 从用户模式切换到主管模式，
并跳转到一个内核指定的入口点。
一旦 CPU 切换到主管模式，
内核就可以验证系统调用的参数（例如，
检查传递给系统调用的地址是否是应用程序内存的一部分），决定是否
允许该应用程序执行所请求的操作（例如，
检查是否允许该应用程序写入指定的文件），然后拒绝它或执行它。
内核控制切换到主管模式的入口点是至关重要的；
如果应用程序可以决定内核入口点，恶意应用程序就可以（例如）在跳过参数验证的地方进入内核。
%% 
\section{Kernel organization}
%% 

一个关键的设计问题是操作系统的哪一部分应该运行在主管模式下。
一种可能性是整个操作系统都驻留在内核中，以便所有系统调用的实现
都运行在主管模式下。
这种组织形式被称为\indextext{monolithic kernel}（宏内核）。

在宏内核组织中，整个操作系统由一个在主管模式下运行的单一程序组成。
这种组织形式很方便的一个原因是操作系统设计者
不需要将代码划分为需要和不需要主管特权的部分。
此外，操作系统的不同部分很容易进行合作，因为它们是一个单一程序的组成部分。
例如，宏内核可能会与文件系统和虚拟内存系统共享一个磁盘块缓存。
缺点是宏内核往往会变得庞大而复杂，以至于没有哪一个开发人员能够理解
代码不同部分之间的所有交互；
这是产生 bug 的温床。内核中的 bug 特别麻烦，因为它可能会导致整个计算机崩溃，
或者导致许多应用程序发生故障，
或者使整个计算机容易受到安全攻击。
\indextext{microkernel}（微内核）旨在减少内核中 bug 的发生率。
其思想是将绝对最少的功能放在内核本身中，以便在主管模式下执行的代码很少，并且
使内核易于理解和进行正确性分析。
操作系统的大部分作为用户级服务器进程运行。
例如，文件系统代码将作为一个服务器进程执行，运行在用户模式而不是主管模式下。
\begin{figure}[t]
\center
\includegraphics[scale=0.5]{fig/mkernel.pdf}
\caption{A microkernel with a file-system server}
\label{fig:mkernel}
\end{figure}

Figure~\ref{fig:mkernel} 阐明了这种微内核设计。在图中，文件系统作为一个
用户级服务器进程运行。
为了允许应用程序与文件服务器交互，内核提供了一种进程间通信机制，用于将消息从一个
用户模式进程发送到另一个进程。
例如，如果像 shell 这样的应用程序想要读取或写入文件，它会向文件服务器发送一条消息并等待
响应。
在微内核中，内核接口由少数几个低级函数组成，用于启动应用程序、发送消息、
访问设备硬件等。这种组织形式允许内核
相对简单，因为大部分操作系统都驻留在用户级服务器中。
在现实世界中，宏内核和微内核都很受欢迎。许多
Unix 内核是宏内核。
例如，Linux 拥有一个宏内核，
尽管一些操作系统功能作为用户级服务器运行（例如窗口系统）。
Linux 为操作系统密集型应用程序提供高性能，部分原因在于内核的子系统可以紧密集成。
像 Minix、L4 和 QNX 这样的操作系统被组织为带有服务器的微内核，并在嵌入式环境中得到了广泛部署。
L4 的一个变体 seL4 非常小，以至于它已经被验证了内存安全性和其他安全属性~\cite{sel4}。
在操作系统的开发人员中，关于哪种组织形式更好存在很多争论，但还没有得出哪一种更好的确凿证据。
此外，这在很大程度上取决于“更好”意味着什么：
更快的性能、更小的代码尺寸、内核的可靠性、
整个操作系统的可靠性（包括用户级服务）等。

还有一些实际的考虑可能比采用哪种组织形式的问题更重要。
一些操作系统具有微内核，但出于性能原因在内核空间运行某些用户级服务。
一些操作系统具有宏内核，因为它们最初就是这样开始的，并且几乎没有动力去转向纯微内核组织，因为新特性可能比重写现有的操作系统以适应微内核设计更重要。
从本书的角度来看，微内核和宏内核操作系统共享许多关键思想。
它们实现系统调用，它们使用页表，它们处理中断，它们支持进程，它们使用锁进行并发控制，它们实现文件系统
等。本书重点关注这些核心思想。
Xv6 是作为一个宏内核实现的，就像大多数 Unix 操作系统一样。
因此，xv6 内核接口对应于操作系统接口，并且内核实现了整个操作系统。
由于 xv6 并没有提供很多服务，其内核比一些微内核还要小，但从概念上讲 xv6 是宏内核。
\section{Code: xv6 organization}

\begin{figure}[t]
\center
\begin{tabular}{l|l|l}
& {\bf File} & {\bf Description} \\
\midrule
Boot & entry.S & Very first boot instructions.
 \\
 & main.c & Control initialization of other modules. \\
 & start.c & Early machine-mode boot code.
 \\
Processes & exec.c & exec() system call. \\
 & proc.c & Processes and scheduling. \\
 & swtch.S & Thread switching.
 \\
 & sysproc.c & Process-related system calls. \\
Traps & kernelvec.S & Handle traps from kernel code.
 \\
 & trampoline.S & Handle traps from user code. \\
 & trap.c & C code to handle and return from traps and interrupts.
 \\
 & syscall.c & Dispatch system calls to handling function. \\
Memory & vm.c & Manage page tables and address spaces.
 \\
 & kalloc.c & Physical page allocator. \\
Devices & console.c & Connect to the user keyboard and screen.
 \\
 & plic.c & RISC-V interrupt controller. \\
 & printf.c & Formatted output to the console.
 \\
 & uart.c & Serial-port console device driver. \\
 & virtio\_disk.c & Disk device driver.
 \\
FS & bio.c & Disk block cache for the file system. \\
 & file.c & File descriptor support.
 \\
 & fs.c & File system. \\
 & log.c & File system logging and crash recovery.
 \\
 & sysfile.c & File-related system calls. \\
 & pipe.c & Pipes.
 \\
Misc & sleeplock.c & Locks that yield the CPU. \\
 & spinlock.c & Locks that don't yield the CPU.
 \\
 & string.c & C string and byte-array library. \\
\end{tabular}
\caption{Xv6 kernel source files.}
\label{fig:source}
\end{figure}

Xv6 内核源代码在 {\tt kernel/} 子目录下。
Figure~\ref{fig:source} 列出了这些文件，分为内核职责的主要领域：启动系统（引导）、创建和控制进程、处理陷阱（中断和系统调用）、分配内存和配置虚拟地址、控制设备以及管理文件系统。
%% 
\section{Process overview}
%% 

在 xv6 中（与其他 Unix 操作系统一样），隔离的单位是一个\indextext{process}（进程）。
进程抽象可以防止一个进程破坏或监视另一个进程的内存、CPU、文件描述符等。它还可以防止进程破坏内核本身，从而使进程无法颠覆内核的隔离机制。
内核必须小心地实现进程抽象，因为一个有 bug 或恶意的应用程序可能会欺骗内核或硬件做出坏事（例如，绕过隔离）。
内核用于实现进程的机制包括用户/主管模式标志、地址空间以及线程的时间片分割。
为了帮助加强隔离，进程抽象为程序提供了一种拥有自己私有机器的幻觉。
进程为程序提供了一个看似私有的内存系统，或\indextext{address space}（地址空间），其他进程无法对其进行读写。
进程还为程序提供了看似属于其自己的 CPU 来执行程序的指令。
Xv6 使用页表（由硬件实现）为每个进程提供其自己的地址空间。
RISC-V 页表将\indextext{virtual address}（虚拟地址，即 RISC-V 指令操作的地址）翻译（或“映射”）为\indextext{physical address}（物理地址，即 CPU 发送到主内存的地址）。
\begin{figure}[t]
\centering
\includegraphics[scale=0.5]{fig/as.pdf}
\caption{Layout of a process's virtual address space}
\label{fig:as}
\end{figure}

Xv6 为每个进程维护一个单独的页表，用以定义该进程的地址空间。
如 Figure~\ref{fig:as} 所示，
一个地址空间包括从虚拟地址零开始的进程的\indextext{user memory}（用户内存）。
指令首先排在最前面，
接着是全局变量，然后是栈，
最后是一个进程可以根据需要扩展的“堆（heap）”区（用于 malloc）。
有许多因素限制了
进程地址空间的最大规模：
RISC-V 上的指针是 64 位宽的；
硬件在页表中查找虚拟地址时仅使用低 39 位；
而 xv6 仅使用这 39 位中的 38 位。
因此，最大地址是 $2^{38}-1$ = 0x3fffffffff，即 \lstinline{MAXVA}~\lineref{kernel/riscv.h:/define.MAXVA/}。
在地址空间的顶部，xv6 放置了一个 \indextext{trampoline}（蹦床）页（4096 字节）和一个 \indextext{trapframe}（陷阱帧）页。
Xv6 使用这两个页来转换进入内核以及返回；
蹦床页包含进入和退出内核的代码，而陷阱帧是内核保存进程的用户寄存器的地方，正如 Chapter~\ref{CH:TRAP} 所解释的那样。
Xv6 内核为每个进程维护许多状态片段，并将它们收集到 \indexcode{struct proc} \lineref{kernel/proc.h:/^struct.proc/} 中。
一个进程最重要的内核状态片段是它的页表、它的内核栈以及它的运行状态。
我们将使用表示法 \indexcode{p->xxx} 来引用 \lstinline{proc} 结构的元素；例如，\indexcode{p->pagetable} 是一个指向进程页表的指针。
此时，请阅读定义了 {\tt struct proc} 的 {\tt kernel/proc.h}。
理解 xv6 的代码对你来说比这本书更重要；
你应该优先考虑代码，并在需要时参考本书以澄清代码。
某些代码的目的起初可能并不明显，但进一步阅读和搜索代码将会有所帮助。
欢迎随意探索和修改代码。

每个进程都有一个控制线程（或者简称为\indextext{thread}（线程）），它保存了执行该进程所需的状态。
在任何给定的时间，一个线程可能正在 CPU 上执行，或者处于挂起状态（没有执行，但能够在未来恢复执行）。
为了在进程之间切换 CPU，内核会挂起当前在该 CPU 上运行的线程并保存其状态，然后恢复另一个进程先前挂起的线程的状态。
线程的大部分状态（局部变量、函数调用返回地址）都存储在线程的栈中。
每个进程都有两个栈：一个用户栈和一个内核栈（\indexcode{p->kstack}）。
当进程正在执行用户指令时，只有它的用户栈在使用中，而它的内核栈是空的。
当进程进入内核时（由于系统调用或中断），内核代码在进程的内核栈上执行；
当一个进程在内核中时，它的用户栈仍然包含保存的数据，但并没有被活跃地使用。
进程的线程在交替地活跃使用它的用户栈和它的内核栈。
内核栈是独立的（并且受到保护而不受用户代码影响），这样即使进程破坏了其用户栈，内核也可以执行。
进程的用户代码可以通过执行 RISC-V \indexcode{ecall} 指令来发起系统调用。
该指令切换到主管模式并改变程序计数器到一个内核定义的入口点。
入口点处的代码切换到进程的内核栈，并执行实现系统调用的内核指令。
当系统调用完成时，内核通过执行 \indexcode{sret} 指令返回到用户空间，该指令会切换到用户模式并紧随系统调用指令之后恢复执行用户指令。
进程的线程可以在内核中“阻塞（block）”以等待 I/O，并在 I/O 完成时从上次中断的地方恢复。
\indexcode{p->state} 指示进程是被分配、准备运行、当前在 CPU 上运行、等待 I/O 还是正在退出。
\indexcode{p->pagetable} 保存进程的页表，格式符合 RISC-V 硬件的预期。
Xv6 在用户空间执行该进程时，会导致分页硬件使用进程的 \lstinline{p->pagetable}。
进程的页表还用作记录分配用于存储进程内存的物理页地址的记录。
总之，一个进程捆绑了两个设计思想：一个给进程带来其自身内存幻觉的地址空间，以及一个给进程带来其自身 CPU 幻觉的线程。
在 xv6 中，一个进程由一个地址空间和一个线程组成。
在真实的操作系统中，一个进程可能拥有多个线程以利用多个 CPU 的优势。
%% 
\section{Code: starting xv6, the first process and system call}
%% 
为了使 xv6 更加具体，我们将概述内核如何启动和运行第一个进程。
后续章节将更详细地描述在此概述中出现的机制。
请阅读 {\tt kernel/entry.S}、{\tt kernel/start.c}、{\tt kernel/main.c} 和 {\tt user/init.c}。
当 RISC-V 计算机上电时，它会初始化自身并运行一个存储在只读内存中的引导加载程序（boot loader）。
引导加载程序将 xv6 内核复制到物理地址 \texttt{0x80000000} 处的内存中。
它将内核放置在 \texttt{0x80000000} 而不是 \texttt{0x0} 的原因是因为地址范围 \texttt{0x0:0x80000000} 包含 I/O 设备。
然后引导加载程序跳转到从 \indexcode{_entry} \lineref{kernel/entry.S:/^.entry:/} 开始的 xv6。
RISC-V 在启动时禁用分页硬件：虚拟地址直接映射到物理地址。
\lstinline{_entry} 处的指令设置了一个栈，以便 xv6 可以运行 C 代码。
Xv6 在文件 \lstinline{start.c} \lineref{kernel/start.c:/stack0/} 中为这个栈 \lstinline{stack0} 声明了空间。
\lstinline{_entry} 处的代码将地址 \lstinline{stack0+4096}（栈的顶部）加载到栈指针寄存器 \texttt{sp} 中，因为 RISC-V 上的栈是向下增长的。
现在内核有了一个栈，\lstinline{_entry} 调用 \lstinline{start} \lineref{kernel/start.c:/^start/} 处的 C 代码。
函数 \lstinline{start} 执行一些 CPU 仅在机器模式下允许的设置，最关键的是对时钟芯片进行编程以产生定时器中断。
然后 {\tt start} 使用 RISC-V 的 {\tt mret} 指令切换到主管模式并跳转到 \lstinline{main} \lineref{kernel/main.c:/^main/}。
{\tt mret} 需要进行一点设置：
{\tt start} 将寄存器 \lstinline{mstatus} 中的前一个特权模式设置为核心态（supervisor），通过将 \lstinline{main} 的地址写入寄存器 \lstinline{mepc} 来将目标地址设置为 \lstinline{main}，通过将 \lstinline{0} 写入页表寄存器 \lstinline{satp} 来在主管模式下禁用虚拟地址转换，并将所有中断和异常委托给主管模式。
在 \lstinline{main} \lineref{kernel/main.c:/^main/} 初始化了几个设备和子系统之后，它通过调用 \lstinline{userinit} \lineref{kernel/proc.c:/^userinit/} 创建了第一个进程。
所有新创建的进程都在内核中的 \lstinline{forkret} \lineref{kernel/proc.c:/^forkret/} 开始执行。
作为第一个进程的特殊情况，{\tt forkret} 调用 {\tt kexec} 来加载用户程序 {\tt /init}。
在调用 \lstinline{kexec} 之后，{\tt forkret} 返回到 {\tt /init} 进程中的用户空间。
\lstinline{init} \lineref{user/init.c:/^main/} 在需要时创建一个新的控制台设备文件，然后将其打开为文件描述符 0、1 和 2。然后它在控制台上启动一个 shell。
系统已经启动。

\section{Security Model}

你可能会想操作系统是如何处理有 bug 的或恶意的代码的。
因为应对恶意严格来说比处理意外的 bug 更困难，所以将主要精力集中在提供防御恶意的安全性上是合理的。
这是操作系统设计中典型安全假设和目标的高级视图。
操作系统必须假设进程的用户级代码将尽其最大努力去破坏内核或其他进程。
用户代码可能会尝试解引用其允许的地址空间之外的指针；它可能会尝试执行不适用于用户代码的指令；
它可能会尝试读取和写入 RISC-V 控制寄存器；它可能会尝试访问设备硬件；
并且它可能会向系统调用传递精明的值，试图诱骗内核崩溃或做出愚蠢的事情。
内核的目标是限制每个用户进程，使其只能访问自己的用户内存，使用 32 个通用 RISC-V 寄存器，并以系统调用预期允许的方式影响内核和其他进程。
内核必须阻止任何其他行为。这些通常是内核设计中的绝对要求。
对内核自身代码的期望是不同的。内核代码被假定是由用意良好且细心的程序员编写的、是没有 bug 的，并且不包含任何恶意内容。
这一假设影响了我们如何分析内核代码。例如，有许多内部内核函数（例如自旋锁），如果内核代码不正确地使用它们，将会导致严重的问题。
然而，我们假设内核正确地使用了它自己的函数。
在硬件层面上，RISC-V CPU、RAM、磁盘等被假定按照文档中的广告进行操作，没有硬件 bug。
现实生活并非如此简单。很难防止滥用的用户程序通过消耗内核保护的资源（磁盘空间、CPU 时间、进程表项等）来以使系统无法使用的方式调用系统调用。
通常不可能编写 100\% 无 bug 的内核代码或设计无 bug 的硬件；
如果恶意用户代码的编写者意识到了内核或硬件的 bug，他们就会利用它们。
即使在成熟、广泛使用的内核（如 Linux）中，人们也经常发现以前未知的漏洞~\cite{mitre:cves}。
最后，用户代码和内核代码之间的区别有时是模糊的：一些特权用户级进程可能会提供基本服务，并且实际上是操作系统的一部分，而在一些操作系统中，特权用户代码可以将新代码插入到内核中（如 Linux 的可加载内核模块和 eBPF）。
作为针对内核 bug 的部分防御，xv6 代码包括了对不一致和不可恢复错误的检查，并将调用 {\tt panic()} 做出“恐慌（panic）”响应。
该函数打印一条错误消息并停止系统。恐慌是不希望看到的，但比继续执行要好。
通常，恐慌是由于内核 bug 导致内核数据不正确或导致内核执行非法操作（例如引用不存在的内存）而引起的；
在这种情况下，使用 {\tt panic()} 停止执行比尝试在不一致的状态下继续运行更安全。
内核开发人员对恐慌的反应将是努力识别并修复底层的代码 bug。
%% 
\section{Real world}
%% 

大多数操作系统都采用了进程概念，而且大多数进程看起来与 xv6 的相似。
然而，现代操作系统在一个进程内支持多个线程，以允许单个进程利用多个 CPU。
在一个进程中支持多个线程涉及到相当多的 xv6 所没有的机制，通常包括接口更改（例如 Linux 的 \lstinline{clone}，它是 \lstinline{fork} 的一种变体），以控制线程共享进程的哪些方面。
%% 
\section{Exercises}
%%

\begin{enumerate}

\item Add a system call to xv6
  that returns the amount of free memory available.
% break *0x3ffffff000
% disas 0x3ffffff000,+8
% set disassemble-next-line auto
% x/i $pc
% break *0x3ffffff10e
% print/x $sepc
% print/x $pc

\end{enumerate}
